\documentclass[10pt,a4paper]{amsart}
\usepackage[margin=19mm]{geometry}
\usepackage{amsmath,amssymb,mathtools}
\usepackage[scr=rsfs,cal=euler]{mathalfa}
\usepackage{xfrac}
\usepackage[unicode,hidelinks,bookmarks=false]{hyperref}
\newcommand{\ie}{i.e.\ }
\newcommand{\cf}{cf.\ }
\newcommand{\resp}{respectively\ }
\newcommand{\Subsection}{\subsection}
\providecommand{\nobreakdash}{\nobreak\hbox{-}}
\setlength{\emergencystretch}{2em}
\multlinegap0pt
\usepackage{bm}
\usepackage{amssymb}
\newtheorem{theorem}{Theorem}
\newcommand{\NN}{{\Bbb N}}
\newcommand{\rd}{\,{\rm d}}
\title{Lebesgue constants for the Walsh system and the discrepancy of the van der Corput sequence}
\author{Josef Dick, Friedrich Pillichshammer}
\date{Source: arXiv:2412.03164v1 (CC BY 4.0)}
\begin{document}
\maketitle

\noindent\emph{Source and licence.} Lebesgue constants for the Walsh system and the discrepancy of the van der Corput sequence. Version arXiv:2412.03164v1 (CC BY 4.0), distributed by arXiv under the Creative Commons Attribution 4.0 International licence, http://creativecommons.org/licenses/by/4.0/, as recorded in the arXiv licence metadata for this exact version and shown on the abstract page https://arxiv.org/abs/2412.03164v1. Full licence notice: \texttt{licenses/arxiv-2412.03164-CC-BY-4.0.txt}.\par\smallskip

\noindent\textit{Editorial scope.} Counted unit: the article's theorem thm2 (source0.tex lines 115--117). The article's complete original proof of it is delivered with it (source0.tex lines 119--134, one \texttt{proof} environment of 16 lines). No mathematics is written, repaired or re-derived: the statement and the proof are the article's own lines, in the article's own notation, and no retained line is substituted or re-flowed.  Retained uncounted as the source-local context the proof needs: lines 84--97; lines 86--88.  Nothing was omitted from any retained span, so the package carries no omission record.  The retained lines cite 2 works, whose entries are transcribed verbatim from the article's own bibliography source in the e-print and delivered with short editorial labels recorded in the item's reference mapping.  No retained line contains a figure environment, an image-inclusion command or drawing code, so no image is delivered and the package declares no asset dependency.  Editorial formatting only, in addition: an AMS \texttt{amsart} preamble that restates the article's own declarations and macros used by the retained lines, namely the environments \texttt{theorem} on separate counters, together with the standard packages those lines need.  The retained environments print in this document as Theorem 1 in order of appearance, whereas the article prints the same items with the numbering of its own full document.  The complete native source of the article as distributed in the arXiv e-print for this exact version is bundled unchanged as \texttt{sources/169597/source0.tex}.
\begin{theorem}[Fine]\label{thm:fine}
The Lebesgue constants $L_n$ of the Walsh system satisfy 
\begin{equation}\label{fo:LCW}
L_n=\nu-\sum_{1 \le j < i \le \nu} 2^{n_i-n_{j}},
\end{equation}
where $n$ is of the form $n=2^{n_1}+2^{n_2}+\cdots+2^{n_{\nu}}$ with integer exponents $n_1 > n_2 > \cdots > n_{\nu}\ge 0$. Furthermore, the following properties hold:
\begin{enumerate}
\item $L_{2n}=L_n$ and $L_{2n+1}=(1+L_n+L_{n+1})/2$.
\item $L_n=O(\log n)$.
\item $\frac{1}{n} \sum_{k=1}^n L_k =\frac{\log n}{4 \log 2}+O(1)$.
\item $\limsup_{n \rightarrow \infty} \left[L_n-\left(\frac{4}{9}+\frac{\log 3}{3 \log 2}+\frac{\log n}{3 \log 2}\right)\right]=0$.
\item the generating function for $L_n$ is $\sum_{n=1}^{\infty} L_n z^n =\frac{1}{2} \frac{z}{(1-z)^2} \sum_{k=0}^{\infty}\frac{1}{2^k} \frac{1-z^{2^k}}{1+z^{2^k}}$ for $|z|<1$.
\end{enumerate}
\end{theorem}

\begin{equation}
L_n=\nu-\sum_{1 \le j < i \le \nu} 2^{n_i-n_{j}},
\end{equation}

\begin{theorem}\label{thm2}
The $n$-th Lebesgue constant $L_n$ for the Walsh system is exactly the nonnormalized star discrepancy $d_n$ of the first $n$ terms of the van der Corput sequence, that is $$d_n=L_n.$$ 
\end{theorem}

\begin{proof}
Let $n \in \NN$ be of the form $n=2^{n_1}+2^{n_2}+\cdots+2^{n_{\nu}}$ with integer exponents $n_1 > n_2 > \cdots > n_{\nu}\ge 0$. For the initial $n$ terms of $Y^{\rm vdC}$ we have 
\begin{align*}
\{y_0,y_1,\ldots,y_{n-1}\} = & \bigcup_{i=1}^{\nu}\left\{y_{\ell}\ : \ \ell \in \{2^{n_1}+\cdots+2^{n_{i-1}},\ldots, 2^{n_1}+\cdots+2^{n_{i-1}}+2^{n_i}-1\}\right\}\\
= & \bigcup_{i=1}^{\nu} \left\{\frac{k}{2^{n_i}}+\frac{1}{2^{n_{i-1}+1}}+\frac{1}{2^{n_{i-2}+1}}+\cdots +\frac{1}{2^{n_1+1}} \ : \ k \in \{0,1,\ldots,2^{n_i}-1\}\right\},
\end{align*}
where we put $2^{n_1}+\cdots+2^{n_{i-1}}:=0$ if $i=1$.

It is well known that the star discrepancy of the van der Corput sequence is nonnegative (see \cite[Remark~3]{p04}) and twice the $L_1([0,1])$-norm of the discrepancy function (see, for example, \cite{PA}). Hence 
\begin{align}\label{stl1}
d_n = & 2 n  \int_0^1 \Delta_{Y^{\rm vdc},n}(t) \rd t = 2 \int_0^1 \sum_{\ell=0}^{n-1} (\boldsymbol{1}_{[0,t)}(y_{\ell}) -t) \rd t\\
= & 2 \sum_{\ell=0}^{n-1} \left(\frac{1}{2}-y_{\ell}\right) = 2 \sum_{i=1}^{\nu} \sum_{k=0}^{2^{n_i}-1} \left(\frac{1}{2}-\frac{k}{2^{n_i}}- \sum_{j=1}^{i-1} \frac{1}{2^{n_j+1}}\right)\nonumber\\
= & 2 \left(\frac{\nu}{2}-\sum_{i=1}^{\nu}2^{n_i} \sum_{j=1}^{i-1} \frac{1}{2^{n_j+1}} \right) = \nu-\sum_{1\le j < i \le \nu}2^{n_i-n_j} =L_n,\nonumber
\end{align}
where the last equality is \eqref{fo:LCW} from Theorem~\ref{thm:fine}.
\end{proof}

\begin{thebibliography}{99}
\bibitem[PA]{PA} P.D.~Pro\u{\i}nov, E.Y.~Atanassov: On the distribution of the van der Corput generalized sequences. C. R. Acad. Sci. Paris S\'er. I Math. 307(18): 895--900, 1988. %
\bibitem[p04]{p04} F.~Pillichshammer: On the discrepancy of $(0,1)$-sequences. J. Number Theory 104: 301--314, 2004. %
\end{thebibliography}
\end{document}
